% Generated draft: requires mathematical and source review.
\documentclass[11pt]{article}
\usepackage[margin=25mm]{geometry}
\usepackage{fontspec}
\setmainfont{Latin Modern Roman}
\usepackage{amsmath,amssymb,mathtools,unicode-math}
\setmathfont{Latin Modern Math}
\usepackage{enumitem,needspace,xurl,hyperref}
\setlength{\parindent}{0pt}
\setlength{\parskip}{5pt}
\setlength{\emergencystretch}{3em}
\newcommand{\N}{\mathbb{N}}
\newcommand{\Z}{\mathbb{Z}}
\newcommand{\Q}{\mathbb{Q}}
\newcommand{\R}{\mathbb{R}}
\newcommand{\C}{\mathbb{C}}
\newcommand{\F}{\mathbb{F}}
\newcommand{\A}{\mathbb{A}}
\newcommand{\PP}{\mathbb{P}}
\newcommand{\E}{\mathbb{E}}
\newcommand{\eps}{\varepsilon}
\newcommand{\dd}{\,\mathrm{d}}
\begin{document}
\section*{UnsolvedMath ID: 30005423}
\textbf{Source problem code:} catalogue entry
\subsection{Definitions and mathematical statement}
Let \(p\) be a prime and put \(R_p=\Z/p^6\Z\). Here \(\operatorname{mod}R_p\) denotes the category of finitely generated \(R_p\)-modules; these are precisely finite abelian groups annihilated by \(p^6\), equipped with their natural \(R_p\)-module structures. The exponent is required to divide \(p^6\), not necessarily to equal \(p^6\).

Define \(\operatorname{mono}(1\to2,\operatorname{mod}R_p)\) to have as objects injective \(R_p\)-linear maps \(f:A\hookrightarrow B\), where \(A,B\in\operatorname{mod}R_p\). A morphism from \(f:A\hookrightarrow B\) to \(g:A'\hookrightarrow B'\) is a pair of \(R_p\)-linear maps \((\alpha:A\to A',\beta:B\to B')\) satisfying
\[
\beta\circ f=g\circ\alpha.
\]
Under the usual interpretation of Birkhoff equivalence, \(f\) and \(g\) are equivalent exactly when such a pair exists with both \(\alpha\) and \(\beta\) isomorphisms. Equivalently, the problem concerns pairs \((B,U)\), with \(B\) a finite abelian group annihilated by \(p^6\) and \(U\leq B\), where \((B,U)\) and \((B',U')\) are equivalent if there is a group isomorphism \(\beta:B\to B'\) with \(\beta(U)=U'\).

\textbf{Classification problem.} Describe the isomorphism classes of all objects of
\[
\operatorname{mono}(1\to2,\operatorname{mod}(\Z/p^6\Z)),
\]
that is, classify all these embeddings up to the equivalence just specified, for arbitrary finite \(A\) and \(B\).

\textbf{Unresolved source qualifications.} The supplied record does not explicitly quantify \(p\); the formulation above treats it as an arbitrary prime, with the classification sought for each such \(p\). The record also does not define Birkhoff equivalence, so the isomorphism-of-embeddings interpretation requires source confirmation. Its original excerpt asks to describe the category and then contains a truncated discussion of general quiver representations and separated monomorphisms. That incomplete discussion does not specify an additional usable hypothesis. The precise desired format of a classification is not supplied.
\subsection{Short English statement}
For each prime \(p\), classify embeddings between finite abelian groups annihilated by \(p^6\), identifying two embeddings when isomorphisms of their domains and codomains make the resulting square commute. The interpretation of Birkhoff equivalence and the scope of the requested classification require human source review. The supplied curation labels the problem open.
\subsection{Sources}
\begin{itemize}
\item Local UnsolvedMath registry: \texttt{data/unsolvedmath/problems.json}, permanent ID \(30005423\), source problem code \texttt{catalogue entry}, in the Oberwolfach Reports Open Problems collection. This draft uses the supplied complete registry record; local file contents were not independently inspected.
\item Supplied primary report citation: \emph{Representation Theory of Quivers and Finite-Dimensional Algebras} (2023), Oberwolfach Reports, pp.\ 397--486, DOI \texttt{10.4171/owr/2023/7}. \url{https://doi.org/10.4171/owr/2023/7}. This is also the supplied literature source labelled ``original OWR report.'' The record gives proposed year 2023 and names no individual proposer.
\item Supplied literature source: Gao--K\"ulshammer--Kvamme--Psaroudakis, PLMS (2024). \url{https://londmathsoc.onlinelibrary.wiley.com/doi/full/10.1112/plms.12640}. No article title or more detailed bibliographic information is supplied.
\item The literature assessment was checked at \texttt{2026-08-22}, according to the supplied dated curation, not independent verification for this draft. It labels the problem open and reports that the \(n=6\) classification remains unresolved, while asserting finite-type results for \(n\leq5\) and controlled wildness for \(n\geq7\). These are preserved as supplied literature claims; their precise hypotheses and time-sensitive status require independent verification.
\item The registry labels the statement \texttt{corrected\_verified} and says it was checked against the official Oberwolfach report, preserving its hypotheses, notation, and open mathematical target. This is a supplied verification claim, not an independent source check performed for this draft.
\end{itemize}
\end{document}
